\section{Einleitung}
	Warum können tonnenschwere Flugzeuge fliegen? Die Antwort auf diese fundamentale Frage der Aerodynamik führt direkt zum Satz von Kutta-Joukowski, einem der bedeutendsten Grundsätze der Strömungsmechanik. Das Theorem stellt eine direkte, mathematische Verknüpfung zwischen der Zirkulation einer idealen Strömung um einen Körper und der daraus resultierenden dynamischen Auftriebskraft her. Entwickelt an der Schwelle zum 20. Jahrhundert durch den deutschen Mathematiker Martin Wilhelm Kutta und den russischen Aerodynamiker Nikolai Joukowski, bildet dieser Ansatz bis heute das theoretische Fundament für das Verständnis von Tragflügelprofilen.